% =====================================================
% 题型：立体几何填空题 (q_tjw_20260605_11_oblique_projection.tex)
% =====================================================
% =====================================================
% 难度：★★ (中档)
% =====================================================
\newcommand{\qSolidObliqueProjectionStem}{%
  如图所示的 $\triangle A'B'C'$ 是水平放置的 $\triangle ABC$ 的斜二测画法的直观图，已知 $A'B' = 4$，$A'C' = B'C' = \sqrt{13}$，则在 $\triangle ABC$ 中，$AC = $ \blank{2cm}.%
}

\newcommand{\qSolidObliqueProjectionFigure}[1][1.2]{%
  \begin{tikzpicture}[scale=#1, >=stealth]
    % =====================================================
    % 左图：原平面图形 \triangle ABC
    % =====================================================
    % 坐标轴
    \draw[->, gray, thin] (-0.5,0) -- (3.0,0) node[below] {\small $x$};
    \draw[->, gray, thin] (0,-0.5) -- (0,4.8) node[left] {\small $y$};
    \node[below left, gray] at (0,0) {\tiny $O$};
    
    % 顶点 (原图比例缩小以适应版面，设 x 轴 1 单位 = 0.5cm，y轴缩放)
    \coordinate (A) at (0.5,0); % 对应 OA = 1 (缩放0.5)
    \coordinate (B) at (2.5,0); % 对应 OB = 5 (缩放0.5)
    \coordinate (C) at (0,4.24); % 对应 OC = 6\sqrt{2} \approx 8.48 (缩放0.5)
    
    \draw[thick] (A) -- (B) -- (C) -- cycle;
    \draw[thick, red] (0,0) -- (A); % OA
    \draw[thick, blue] (0,0) -- (C); % OC
    
    \fill (0,0) circle (1.2pt);
    \fill (A) circle (1.2pt);
    \fill (B) circle (1.2pt);
    \fill (C) circle (1.2pt);
    
    \node[below] at (A) {\small $A$};
    \node[below] at (B) {\small $B$};
    \node[left] at (C) {\small $C$};
    
    \node[below, red] at (0.25,0) {\small $1$};
    \node[left, blue] at (0, 2.12) {\small $6\sqrt{2}$};
    
    % 直角符号
    \draw[gray] (0.15, 0) -- (0.15, 0.15) -- (0, 0.15);
    
    \node[below] at (1.0, -0.6) {\textbf{原平面图形 $\triangle ABC$}};
    
    % =====================================================
    % 右图：斜二测直观图 \triangle A'B'C'
    % =====================================================
    \begin{scope}[xshift=5.2cm]
      % 坐标轴
      \draw[->, gray, thin] (-0.5,0) -- (3.0,0) node[below] {\small $x'$};
      \draw[->, gray, thin] (0,-0.5) -- (0,2.5) node[left] {\small $y'$ (垂直)};
      \draw[->, blue, thin] (0,0) -- (45:3.0) node[above right] {\small $y'$ (斜轴)};
      \node[below left, gray] at (0,0) {\tiny $O'$};
      
      % 顶点 (比例同样为 0.5)
      \coordinate (Oprime) at (0,0);
      \coordinate (Aprime) at (0.5,0); % 对应 O'A' = 1
      \coordinate (Bprime) at (2.5,0); % 对应 O'B' = 5，即 A'B' = 4
      % C' 在 y' 轴上，O'C' = 3\sqrt{2} \approx 4.24
      % 页面坐标 C' 为 (4.24 * 0.5 * cos 45, 4.24 * 0.5 * sin 45) = (1.5, 1.5)
      \coordinate (Cprime) at (1.5, 1.5); 
      % C' 向 x' 轴作垂线
      \coordinate (Dprime) at (1.5, 0); 
      
      \draw[thick] (Aprime) -- (Bprime) -- (Cprime) -- cycle;
      
      % 线段 O'A' 和 O'C'
      \draw[thick, red] (Oprime) -- (Aprime);
      \draw[thick, blue] (Oprime) -- (Cprime);
      % 垂直高度 C'D'
      \draw[thick, dashed, gray] (Cprime) -- (Dprime);
      
      \fill (Oprime) circle (1.2pt);
      \fill (Aprime) circle (1.2pt);
      \fill (Bprime) circle (1.2pt);
      \fill (Cprime) circle (1.2pt);
      \fill (Dprime) circle (1.2pt);
      
      \node[below] at (Aprime) {\small $A'$};
      \node[below] at (Bprime) {\small $B'$};
      \node[above left] at (Cprime) {\small $C'$};
      \node[below] at (Dprime) {\small $D'$};
      
      \node[below, red] at (0.25, 0) {\small $1$};
      \node[above left, blue] at ($(Oprime)!0.5!(Cprime)$) {\small $3\sqrt{2}$};
      \node[right, gray] at ($(Cprime)!0.5!(Dprime)$) {\small $3$};
      \node[below, gray] at ($(Aprime)!0.5!(Bprime)$) {\small $4$};
      
      % 直角符号 (D')
      \draw[gray] (1.35, 0) -- (1.35, 0.15) -- (1.5, 0.15);
      
      \node[below] at (1.25, -0.6) {\textbf{直观图 $\triangle A'B'C'$}};
    \end{scope}
  \end{tikzpicture}%
}

\newcommand{\qSolidObliqueProjectionAnswer}{$\sqrt{73}$}

\newcommand{\qSolidObliqueProjectionSolution}{%
  【解题思路】\\
  这道题最核心的“题眼”在于观察直观图中点与坐标轴的位置关系。根据直观图，点 \(C'\) 在斜轴 \(y'\) 上，点 \(A', B'\) 在轴 \(x'\) 上。由于直角坐标系下 \(x\) 轴与 \(y\) 轴天然垂直，我们可以在直观图求出各线段长度后，直接在原图中利用直角三角形的勾股定理求出 \(AC\)，从而避免繁琐的坐标计算。\\
  【解析计算】\\
  第一步：在直观图 \(\triangle A'B'C'\) 中，找出原点 \(O'\) 相关的线段长度。\\
  根据原图图像，点 \(C'\) 落在斜 \(y'\) 轴上，点 \(A', B'\) 落在 \(x'\) 轴上。\\
  在等腰 \(\triangle A'B'C'\) 中，作 \(C'D' \perp x'\) 轴于点 \(D'\)。\\
  已知 \(A'B' = 4\)，\(A'C' = B'C' = \sqrt{13}\)，由等腰三角形三线合一得 \(D'\) 为 \(A'B'\) 中点，且 \(A'D' = 2\)。\\
  在直角 \(\triangle A'D'C'\) 中，高 \(C'D' = \sqrt{(\sqrt{13})^2 - 2^2} = 3\)。\\
  因为点 \(C'\) 在斜轴 \(y'\) 上（倾角 \(45^\circ\)），在直角 \(\triangle O'D'C'\) 中，\(\angle C'O'D' = 45^\circ\)。\\
  所以 \(\triangle O'D'C'\) 是等腰直角三角形，\\
  这就意味着水平偏移量 \(O'D' = C'D' = 3\)。\\
  由于 \(\angle C'O'D' = 45^\circ\)，斜轴上的线段长度为 \(O'C' = \dfrac{C'D'}{\sin 45^\circ} = 3\sqrt{2}\)。\\
  因为点 \(A'\) 在点 \(D'\) 左侧且距离为 \(2\)，所以线段 \(O'A' = O'D' - A'D' = 3 - 2 = 1\)。\\
  第二步：利用斜二测规则还原到原平面图形 \(\triangle ABC\) 中。\\
  根据“横不变，竖减半，平行关系不变”原则：\\
  1. \(A\)、\(C\) 还原后必然分别在 \(x\) 轴和 \(y\) 轴上。且长度翻倍与保留关系为：\(OA = O'A' = 1\)，\(OC = 2O'C' = 6\sqrt{2}\)。\\
  2. 原直角坐标系中，\(x\) 轴与 \(y\) 轴天然垂直。\\
  所以，原点 \(O\)、点 \(C\)、点 \(A\) 构成的 \(\triangle COA\) 必然是一个直角三角形（\(\angle COA = 90^\circ\)）！\\
  第三步：在 Rt\(\triangle COA\) 中直接求 \(AC\)。\\
  利用最简单的勾股定理：\\
  \(AC = \sqrt{OA^2 + OC^2} = \sqrt{1^2 + (6\sqrt{2})^2} = \sqrt{1 + 72} = \sqrt{73}\)。%
}

\newcommand{\qSolidObliqueProjectionDiagnosis}{%
  本题学生极其容易被“设坐标代入距离公式”的代数计算思维束缚，忽略了仔细观察图像中 $C'$ 恰好在 $y'$ 轴上这一得天独厚的几何条件。一旦强行设未知点坐标，不仅计算量庞大，而且极易在等式推导中全盘崩溃。
}

\newcommand{\qSolidObliqueProjectionTeachingNotes}{%
  \item \textbf{强调“看图”的几何直觉}：
    \begin{itemize}
       \item 带着学生看原卷配图，指出 $C'$ 在 $y'$ 轴上这最关键的一点。
       \item 启发提问：“如果在直观图里，两个点分别在 $x'$ 和 $y'$ 两根轴上，还原回原图，它们连成什么三角形？”
       \item 总结方法论：\textbf{斜二测还原问题，凡是落在了轴上的点，必产生垂直关系！优先用直角三角形的勾股定理去算，绝对不要去硬算坐标！}这不仅计算量极小而且完全不会出错。
    \end{itemize}
}


% =====================================================
% 别名兼容：旧课次宏定义
% =====================================================
\newcommand{\qTJWObliqueProjectionStem}{\qSolidObliqueProjectionStem}
\newcommand{\qTJWObliqueProjectionFigure}[1][1.2]{\qSolidObliqueProjectionFigure[#1]}
\newcommand{\qTJWObliqueProjectionAnswer}{\qSolidObliqueProjectionAnswer}
\newcommand{\qTJWObliqueProjectionSolution}{\qSolidObliqueProjectionSolution}
\newcommand{\qTJWObliqueProjectionDiagnosis}{\qSolidObliqueProjectionDiagnosis}
\newcommand{\qTJWObliqueProjectionTeachingNotes}{\qSolidObliqueProjectionTeachingNotes}
